
\subsubsection{完整工程附录}
\label{sec:ros2-project}
\begin{example}
在工作空间 src 下创建 my\_package 与 my\_msgs 两个包。按下列路径保存文件，先编译消息包，再运行节点。
代码中的维护者信息为教学占位值，发布自己的包时应改为实际信息。
\begin{tbash}
mkdir -p ~/ros2_ws/src/my_package/include/my_package
mkdir -p ~/ros2_ws/src/my_package/src ~/ros2_ws/src/my_msgs/msg
\end{tbash}
my\_msgs/msg/My.msg 使用正文的四字段定义；其余文件完整列在下方。


\begin{lstlisting}[language=cmake,label=lst:ament_cmake_auto_config,caption={my\_package/CMakeLists.txt}]
cmake_minimum_required(VERSION 3.8)
project(my_package)
find_package(ament_cmake REQUIRED)
find_package(rclcpp REQUIRED)
find_package(std_msgs REQUIRED)
find_package(my_msgs REQUIRED)
add_executable(my_node src/my_node.cpp src/node_run.cpp)
target_include_directories(my_node PRIVATE include)
target_compile_features(my_node PRIVATE cxx_std_17)
ament_target_dependencies(my_node rclcpp std_msgs my_msgs)
install(TARGETS my_node DESTINATION lib/${PROJECT_NAME})
ament_package()
\end{lstlisting}


\begin{lstlisting}[language=XML,label=lst:xml_1,caption={my\_package/package.xml}]
<?xml version="1.0"?>
<package format="3">
  <name>my_package</name>
  <version>0.0.1</version>
  <description>字符串与自定义消息节点示例</description>
  <maintainer email="student@example.com">student</maintainer>
  <license>Apache-2.0</license>
  <buildtool_depend>ament_cmake</buildtool_depend>
  <depend>rclcpp</depend>
  <depend>std_msgs</depend>
  <depend>my_msgs</depend>
  <export><build_type>ament_cmake</build_type></export>
</package>
\end{lstlisting}

\begin{lstlisting}[language=c++,label=lst:my_node_cpp,caption=my\_node.cpp,numbers=left]
    
    #include "my_package/my_node.h"
    
    #include <chrono>
    #include <functional>
    #include <memory>
    #include <string>
    
    my_node::my_node() : my_node(rclcpp::NodeOptions())
    {
    }
    
    
    my_node::my_node(const rclcpp::NodeOptions& options): Node("my_node", options)
    {
        RCLCPP_INFO(get_logger(), "init");
        string_subscription_ = create_subscription<std_msgs::msg::String>(
            "receive",
            10,
            std::bind(&my_node::subscription_callback, this, std::placeholders::_1)
        );
        string_publisher_ = create_publisher<std_msgs::msg::String>(
            "send",
            10
        );
        timer_publisher_ = create_publisher<std_msgs::msg::String>(
            "timer_",
            10
        );
        timer_ = create_wall_timer(
            std::chrono::milliseconds(1000), // 定时器间隔
            std::bind(&my_node::timer_callback, this)); // 回调函数
    
        my_message_publisher_ = create_publisher<my_msgs::msg::My>(
            "my_msg",
            10
        );
    
        my_message_subscription_ = create_subscription<my_msgs::msg::My>(
            "my_msg",
            10,
            std::bind(&my_node::my_message_callback,this,std::placeholders::_1)
        );
    
    }
    
    void my_node::subscription_callback(const std_msgs::msg::String::SharedPtr msg)
    {
        auto message = std_msgs::msg::String();
        message.data = "I can hear:" + msg->data;
        RCLCPP_INFO(get_logger(), "I hear %s", msg->data.c_str());
        string_publisher_->publish(message);
    
    }
    
    void my_node::timer_callback()
    {
        RCLCPP_INFO(get_logger(), "%d s", ++cnt_);
        auto time_msg = std_msgs::msg::String();
        time_msg.data = "Now is " + std::to_string(cnt_);
        timer_publisher_->publish(time_msg);
        auto self_msg = my_msgs::msg::My();
        self_msg.header.stamp = now();
        self_msg.header.frame_id = "camera_link";
        self_msg.id = 1;
        self_msg.yaw = 0.25; // 弧度
        self_msg.detect_success = true;
        my_message_publisher_->publish(self_msg);
    
    }
    void my_node::my_message_callback(const my_msgs::msg::My::SharedPtr msg)
    {
        RCLCPP_INFO(get_logger(),"id %d , yaw %f , detect_success %d ",msg->id,msg->yaw,msg->detect_success);
    }
    
    
\end{lstlisting}

\begin{lstlisting}[language=c++,label=lst:my_node_h,caption=my\_node.h,numbers=left]
    
    #ifndef MY_NODE_H
    #define MY_NODE_H
    
    #include "rclcpp/rclcpp.hpp"
    #include "std_msgs/msg/string.hpp"
    
    #include "my_msgs/msg/my.hpp"
    
    class my_node: public rclcpp::Node {
        rclcpp::Publisher<std_msgs::msg::String>::SharedPtr string_publisher_;
        rclcpp::Publisher<std_msgs::msg::String>::SharedPtr timer_publisher_;
        rclcpp::Subscription<std_msgs::msg::String>::SharedPtr string_subscription_;
        rclcpp::TimerBase::SharedPtr timer_;
        rclcpp::Publisher<my_msgs::msg::My>::SharedPtr  my_message_publisher_;
        rclcpp::Subscription<my_msgs::msg::My>::SharedPtr my_message_subscription_;
        int cnt_ = 0;
    public:
        my_node() ;
        my_node(const rclcpp::NodeOptions &options);
        void subscription_callback(std_msgs::msg::String::SharedPtr msg);
        void timer_callback();
        void my_message_callback(my_msgs::msg::My::SharedPtr msg);
    };
    
    
    
    #endif //MY_NODE_H
    
\end{lstlisting}

\begin{lstlisting}[language=c++,label=lst:node_run_cpp,caption=node\_run.cpp,numbers=left]


    #include <rclcpp/rclcpp.hpp>
    #include <my_package/my_node.h>

    int main(int argc, char **argv)
    {
        rclcpp::init(argc,argv);
        rclcpp::spin(std::make_shared<my_node>());
        rclcpp::shutdown();
    }

\end{lstlisting}


\begin{lstlisting}[language=cmake,label=lst:cmake_2,caption={my\_msgs/CMakeLists.txt}]
cmake_minimum_required(VERSION 3.8)
project(my_msgs)
find_package(ament_cmake REQUIRED)
find_package(rosidl_default_generators REQUIRED)
find_package(std_msgs REQUIRED)
rosidl_generate_interfaces(${PROJECT_NAME}
  "msg/My.msg" DEPENDENCIES std_msgs)
ament_export_dependencies(rosidl_default_runtime)
ament_package()
\end{lstlisting}

\begin{lstlisting}[language=XML,label=lst:xml_2,caption={my\_msgs/package.xml}]
<?xml version="1.0"?>
<package format="3">
  <name>my_msgs</name>
  <version>0.0.1</version>
  <description>教学消息接口</description>
  <maintainer email="student@example.com">student</maintainer>
  <license>Apache-2.0</license>
  <buildtool_depend>ament_cmake</buildtool_depend>
  <buildtool_depend>rosidl_default_generators</buildtool_depend>
  <depend>std_msgs</depend>
  <exec_depend>rosidl_default_runtime</exec_depend>
  <member_of_group>rosidl_interface_packages</member_of_group>
  <export><build_type>ament_cmake</build_type></export>
</package>
\end{lstlisting}

编译及运行（终端一）：
\begin{tbash}
source /opt/ros/humble/setup.bash
cd ~/ros2_ws
colcon build --packages-up-to my_package
source install/setup.bash
ros2 run my_package my_node
\end{tbash}
终端二加载相同环境后观察自定义消息：
\begin{tbash}
source /opt/ros/humble/setup.bash
source ~/ros2_ws/install/setup.bash
ros2 topic echo /my_msg
\end{tbash}
应约每秒收到 id 为 1、yaw 为 0.25 的消息。按 Ctrl+C 结束观察后，可执行
\begin{tbash}
ros2 topic pub --once /receive std_msgs/msg/String "{data: hello}"
\end{tbash}
终端一应打印收到 hello；另开终端订阅 /send 可看到加上前缀的回复。
\end{example}
